\documentclass[12pt,a4paper]{article}

\usepackage[T2A]{fontenc}
\usepackage[utf8]{inputenc}
\usepackage[russian]{babel}

\usepackage[margin=2.5cm]{geometry}

\usepackage{array}
\usepackage{graphicx}

\setlength{\parindent}{0pt}
\setlength{\parskip}{6pt}

\begin{document}

\begin{center}
    {\LARGE\bfseries Отчёт по экспериментам: сравнение производительности свёртки}\\[6pt]
    \textit{Сравнение собственной реализации двумерной свёртки с реализацией из библиотеки OpenCV (cv::filter2D).}
\end{center}

\section*{1. Цель}

Сравнить время работы собственной реализации свёртки RGB-изображения с реализацией из OpenCV и оценить точность измерений по правилам обработки результатов измерений.

\section*{2. Постановка эксперимента}

Измерялось время выполнения одной и той же операции - свёртки с ядром box blur 3*3 и режимом обработки краёв reflect101. Условия эксперимента:

\begin{itemize}
    \item Изображение случайное, фиксированного размера 2560*1920 пикселей, RGB. Случайное изображение генерируется в самой программе, поэтому не требуется набор реальных картинок; на время работы влияет только размер изображения, а не его содержимое.
    \item Для каждой реализации проводилось 40 измерений и дополнительно 3 «прогревочных» запуска, результаты которых не записывались (первые запуски медленнее из-за «холодного» кеша процессора).
    \item Замерялось время только самой свёртки, без чтения и записи файлов. Использовался монотонный таймер (CLOCK\_MONOTONIC в C и steady\_clock в C++).
    \item OpenCV была ограничена одним потоком (cv::setNumThreads(1)), так как собственная реализация однопоточная.
    \item Эксперименты проводились в среде WSL2 (Ubuntu поверх Windows).
\end{itemize}

\section*{3. Методика обработки результатов}

Полученные серии измерений обрабатывались по следующим шагам:

\begin{itemize}
    \item Построение гистограммы для визуальной проверки распределения и поиска выбросов.
    \item Проверка на нормальность двумя тестами - D'Agostino-Pearson (normaltest) и Шапиро-Уилка (shapiro). Признаком прохождения считается p-value > 0.05 хотя бы в одном тесте.
    \item Вычисление среднего и стандартного отклонения (с поправкой ddof=1). Стандартное отклонение не должно превышать 5-10\% от среднего.
    \item Вычисление доверительного интервала с уровнем доверия 95\% по распределению Стьюдента.
    \item Округление: погрешность - до одной значащей цифры (или до двух, если первая цифра - единица), среднее - до той же точности.
    \item Так как времена двух реализаций разного порядка, итоговое сравнение проведено через относительные погрешности.
\end{itemize}

Погрешность измерительного инструмента (разрешение таймера) составляет порядка наносекунд, что пренебрежимо мало по сравнению с разбросом измерений в миллисекундах. Поэтому в качестве итоговой погрешности взята случайная погрешность (доверительный интервал).

\section*{4. Результаты}

В серии измерений собственной реализации присутствовал один одиночный выброс (около 1024 мс при среднем около 967 мс), из-за которого тесты на нормальность не проходили. Согласно методике, одиночный выброс допускается отбросить; после удаления одного значения (осталось 39 измерений) тесты на нормальность стали проходить. Итоговые данные приведены в таблице.

\textbf{1. Сводка результатов измерений}

\begin{center}
\begin{tabular}{|p{6.2cm}|p{5cm}|p{3.2cm}|}
\hline
\textbf{Показатель} & \textbf{Моя реализация} & \textbf{OpenCV} \\
\hline
Число измерений & 39 (после удаления одного выброса) & 40 \\
\hline
Среднее время, мс & 966.85 & 18.66 \\
\hline
Стандартное отклонение, мс & 13.46 & 0.62 \\
\hline
Отклонение от среднего, \% & 1.39 & 3.30 \\
\hline
p-value (normaltest) & 0.30 & 0.60 \\
\hline
p-value (Shapiro) & 0.15 & 0.19 \\
\hline
Доверительный интервал (95\%), мс & +/-4.36 & +/-0.20 \\
\hline
Итог (с округлением), мс & \textbf{967 +/- 4} & \textbf{18.66 +/- 0.2} \\
\hline
\end{tabular}
\end{center}

\textbf{2. Сравнение реализаций}

\begin{center}
\begin{tabular}{|p{10cm}|p{4.5cm}|}
\hline
\textbf{Сравнение} & \textbf{Значение} \\
\hline
Отношение среднего времени (моя реализация / OpenCV) & 51.8 +/- 0.6 \\
\hline
Во сколько раз OpenCV быстрее & \textbf{примерно в 52 раза} \\
\hline
\end{tabular}
\end{center}

Гистограммы обеих серий прилагаются отдельными файлами (histogram\_my\_implementation.png, histogram\_opencv.png).

\section*{5. Выводы}

\begin{itemize}
    \item OpenCV выполняет ту же свёртку примерно в 52 раза быстрее собственной реализации (51.8 +/- 0.6). Это ожидаемо: OpenCV использует векторные инструкции процессора и другие оптимизации.
    \item Стандартное отклонение составило 1.39\% и 3.30\% от среднего - это укладывается в допустимые 5-10\%, то есть измерения стабильны.
    \item После удаления одного одиночного выброса обе серии прошли тесты на нормальность, поэтому доверительные интервалы для средних корректны.
    \item Итоговые времена: собственная реализация - 967 +/- 4 мс, OpenCV - 18.66 +/- 0.2 мс.
\end{itemize}

\end{document}
